\documentclass[UTF8]{ctexart}
\usepackage{xeCJK}
\usepackage{geometry}
\geometry{left=2cm,right=2cm,top=2.5cm,bottom=2.5cm}
\setlength{\headheight}{12.7pt}

\usepackage{titlesec}
\titleformat{\section}
  {\normalfont\large\bfseries}
  {\thesection}
  {1em}
  {}
\titlespacing*{\section}
  {0pt}
  {3.5ex plus 1ex minus .2ex}
  {2.3ex plus .2ex}

\usepackage{amsmath}
\usepackage{amssymb}
\usepackage{amsthm}
\usepackage{enumitem}
\setlist[enumerate]{itemsep=0pt,topsep=0pt,parsep=0pt,leftmargin=0pt,itemindent=2.5em}
\setlist[itemize]{itemsep=0pt,topsep=0pt,parsep=0pt,leftmargin=0pt,itemindent=2.5em}
\usepackage{fancyhdr}
\pagestyle{fancy}
\fancyhf{}
\usepackage{graphicx}
\usepackage{hyperref}
\usepackage{indentfirst}
\usepackage{lmodern}


\begin{document}
\setlength{\abovedisplayskip}{1pt}
\setlength{\belowdisplayskip}{1pt}

\section{有序求和}

\hypertarget{sec:定义1.1}{}
\noindent\textbf{定义1.1 (有序求和)}

给定\href{../../05 抽象代数/01 群论/01 群的基本概念/01 半群与群#nameddest=sec:定义1.1}{半群}
$(A,+)$, 对任意的自然数$m,k$和映射$f:[\,m,m+k\,]\to A$, 对$j\leqslant k$递归定义
\vspace{3pt}
\[\sum_{i=m}^{m+j}f(i)=\left\{\begin{aligned}
    &f(m),&&j=0,\\
    &\sum_{i=m}^{m+j-1}f(i)+f(m+j),&&0<j\leqslant k.
\end{aligned}\right.\]

\vspace{15pt}

一个序列的求和可以拆成两段子序列的求和之和.

\vspace{10pt}

\hypertarget{sec:定理1.1}{}
\noindent\textbf{定理1.1}

任给半群$(A,+)$, 自然数$m,p,k$ ($p<k$)和映射$f:[\,m,m+k\,]\to A$, 有
\vspace{3pt}
\[
    \sum_{i=m}^{m+k}f(i)=\sum_{i=m}^{m+p}f(i)+\sum_{i=m+p+1}^{m+k}f(i).
\]

\noindent{\kaishu 证明.}

对$k>p$归纳证明. 首先, 当$k=p+1$时依定义即得
\vspace{3pt}
\[
    \sum_{i=m}^{m+p+1}f(i)=\sum_{i=m}^{m+p}f(i)+f(m+p+1);\\[3pt]
\]
其次, 如果命题对$k$成立, 那么对$k+1$有
\vspace{3pt}
\begin{equation*}\begin{aligned}
    \sum_{i=m}^{m+k+1}f(i)&=\sum_{i=m}^{m+k}f(i)+f(m+k+1)\\
    &=\sum_{i=m}^{m+p}f(i)+\sum_{i=m+p+1}^{m+k}f(i)+f(m+k+1)=
    \sum_{i=m}^{m+p}f(i)+\sum_{i=m+p+1}^{m+k+1}f(i).
\end{aligned}\tag*{\qedsymbol}\end{equation*}

\vspace{15pt}

如果两个序列逐位相同, 那么求和相等.

\vspace{10pt}

\hypertarget{sec:定理1.2}{}
\noindent\textbf{定理1.2}

任给半群$(A,+)$, 自然数$m,n,k$和映射$f:[\,m,m+k\,]\to A$,
$g:[\,n,n+k\,]\to A$. 如果
\[
    \forall p\leqslant k:\quad f(m+p)=g(n+p),
\]
那么$\displaystyle\sum_{i=m}^{m+k}f(i)=\sum_{j=n}^{n+k}g(j)$.

\vspace{3pt}
\noindent{\kaishu 证明.}

对$p\leqslant k$归纳证明$\displaystyle\sum_{i=m}^{m+p}f(i)=\sum_{j=n}^{n+p}g(j)$.

首先显然$\displaystyle\sum_{i=m}^{m}f(i)=f(m)=g(n)=\sum_{j=n}^{n}g(j)$;
其次当$p+1\leqslant k$时,
假设$\displaystyle\sum_{i=m}^{m+p}f(i)=\sum_{j=n}^{n+p}g(j)$, 则
\vspace{3pt}
\[
    \sum_{i=m}^{m+p+1}f(i)=\sum_{i=m}^{m+p}f(i)+f(m+p+1)=
    \sum_{j=n}^{n+p}g(j)+g(n+p+1)=\sum_{j=n}^{n+p+1}g(j).\\[3pt]
\]

于是归纳成立, 取$p=k$即证. \hfill $\qedsymbol$

\end{document}